\documentclass[11pt,a4paper]{article}

\usepackage[T1]{fontenc}
\usepackage[utf8]{inputenc}
\usepackage{lmodern}
\usepackage{amssymb,amsmath,amsthm,mathtools}
\usepackage[margin=2.5cm]{geometry}
\usepackage{parskip}
\usepackage[hidelinks]{hyperref}
\usepackage{microtype}

% Theorem environments (arXiv / journal style)
\theoremstyle{plain}
\newtheorem{theorem}{Theorem}[section]
\newtheorem{proposition}[theorem]{Proposition}
\newtheorem{lemma}[theorem]{Lemma}
\newtheorem{corollary}[theorem]{Corollary}

\theoremstyle{definition}
\newtheorem{definition}[theorem]{Definition}
\newtheorem{example}[theorem]{Example}
\newtheorem{remark}[theorem]{Remark}

% Notation
\newcommand{\Ext}{\mathrm{Ext}}
\newcommand{\Shape}{\mathrm{Shape}}
\newcommand{\Pos}{\mathrm{Pos}}
\newcommand{\Type}{\mathbf{Type}}
\newcommand{\id}{\mathrm{id}}
\newcommand{\rootpos}{\mathsf{o}}
\newcommand{\sub}{\mathbin{\downarrow}}
\newcommand{\shift}{\oplus}

\title{Directed Containers and the Induced Comonad}
\author{Extracted from a machine-checked development}
\date{}

\begin{document}
\maketitle

\begin{abstract}
We equip containers with a directed structure---a root position, a
sub-shape operation, and a shift combining positions along sub-shapes---subject
to five equational laws (D1--D5). The extension of a directed container carries
a comonad structure whose counit reads the label at the root and whose
comultiplication re-indexes labels along sub-shapes and shifts. Each of the
three comonad laws is deduced from a corresponding directed-container law (or
pair of laws). The proofs of all stated results have been machine-checked; we
record only the mathematical content.
\end{abstract}

\tableofcontents
\bigskip

%======================================================================
\section{Introduction}
%======================================================================

A container $(S,P)$ determines a polynomial endofunctor
$\Ext\,C\,X = \sum_{s\,:\,S}\,(P(s)\to X)$ on types. Many such functors arising
in practice---streams, non-empty lists, tree zippers, the store comonad---are
not merely functors but \emph{comonads}. The combinatorial data needed to equip
$\Ext\,C$ with a comonad structure can be isolated at the level of the container
itself: this is the notion of a \emph{directed container}, introduced by Ahman,
Chapman, and Uustalu.

Concretely, a directed container adds to $(S,P)$ a distinguished root position
$\rootpos(s)\in P(s)$, a sub-shape $s\sub p\in S$ sitting under each position
$p$, and a shift $p\shift q\in P(s)$ combining a position of $s$ with a
position of the sub-shape. Five laws govern how roots, sub-shapes, and shifts
interact; two of them involve transport of positions along shape equalities,
because positions live in dependent fibres.

The counit and comultiplication on $\Ext\,C$ are then written purely in terms
of this combinatorial data, and the three comonad identities follow one-for-one
from the directed laws. The dependent bookkeeping is packaged once and for all
in an extensionality principle for elements of an extension.

\paragraph{What this note claims.}
In ordinary mathematical language: if a container carries a root, a sub-shape
operation, and a shift satisfying D1--D5, then its extension is a comonad on
types, with
\[
  \varepsilon\langle s,v\rangle \;=\; v\bigl(\rootpos(s)\bigr),
  \qquad
  \delta\langle s,v\rangle \;=\;
  \Bigl\langle s,\;
    p\mapsto\bigl\langle s\sub p,\;
      q\mapsto v(p\shift q)\bigr\rangle
  \Bigr\rangle,
\]
and conversely each comonad law is the image of a directed-container law under
this translation. There is no further data.

%======================================================================
\section{Preliminaries}
%======================================================================

We work in a single universe $\Type$ of types. A \textbf{container} $C$
consists of a type $\Shape(C)$ of shapes and, for each shape $s$, a type
$\Pos(C)(s)$ of positions. The \textbf{extension} at a type $X$ is
\begin{equation}
  \Ext\,C\,X \;=\; \sum_{s\,:\,\Shape(C)}\bigl(\Pos(C)(s)\to X\bigr).
\end{equation}
We write elements as pairs $\langle s,\,v\rangle$. The functorial action on a
function $f:X\to Y$ is
$\Ext\,C\,f\,\langle s,v\rangle = \langle s,\, f\circ v\rangle$.

Because several directed-container laws equate shapes and then require
positions in different fibres to be compared, we record an extensionality
principle for elements of an extension.

\begin{lemma}[Dependent extensionality for extensions]\label{lem:ext-eq}
Let $C$ be a container, $X$ a type, and $s_1,s_2\in\Shape(C)$ with
$h:s_1=s_2$. Given labellings $v_1:\Pos(C)(s_1)\to X$ and
$v_2:\Pos(C)(s_2)\to X$, if
\[
  \forall\, q\in\Pos(C)(s_1).\;
  v_1(q) \;=\; v_2\bigl(h\cdot q\bigr)
\]
(where $h\cdot q$ denotes transport of the position $q$ along $h$), then
\[
  \langle s_1,\,v_1\rangle \;=\; \langle s_2,\,v_2\rangle
  \qquad\text{in }\;\Ext\,C\,X.
\]
\end{lemma}

\begin{proof}
By lean.
\end{proof}

%======================================================================
\section{Directed containers}
%======================================================================

\begin{definition}[Directed container]\label{def:directed}
A \textbf{directed container} $D$ consists of an underlying container
$C=(\Shape(D),\Pos(D))$ together with three operations
\begin{enumerate}
  \item a \textbf{root} $\rootpos:\forall s.\;\Pos(D)(s)$,
  \item a \textbf{sub-shape} $({\sub}):\forall s.\;\Pos(D)(s)\to\Shape(D)$,
  \item a \textbf{shift}
        $({\shift}):\forall s.\;\forall p\in\Pos(D)(s).\;
        \Pos(D)(s\sub p)\to\Pos(D)(s)$,
\end{enumerate}
subject to the five laws below. We write $S$ and $P$ for the shape and
position components when no confusion arises.
\end{definition}

\begin{definition}[Directed-container laws]\label{def:d-laws}
Let $D$ be as in Definition~\ref{def:directed}. The laws are:
\begin{enumerate}
  \item[\textbf{D1.}] Sub-shape under the root is the shape itself:
        \[
          \forall s.\; s\sub\rootpos(s) \;=\; s.
        \]
  \item[\textbf{D4.}] Sub-shapes nest along the shift:
        \[
          \forall s.\;\forall p\in P(s).\;\forall q\in P(s\sub p).\;
          s\sub(p\shift q) \;=\; (s\sub p)\sub q.
        \]
  \item[\textbf{D2.}] Shifting by the root is the identity, transported along
        D1:
        \[
          \forall s.\;\forall q\in P\bigl(s\sub\rootpos(s)\bigr).\;
          \rootpos(s)\shift q \;=\; (\mathrm{D1}_s)\cdot q.
        \]
  \item[\textbf{D3.}] Shifting by the root of the sub-shape is the identity:
        \[
          \forall s.\;\forall p\in P(s).\;
          p\shift\rootpos(s\sub p) \;=\; p.
        \]
  \item[\textbf{D5.}] The shift is associative, with the right argument
        transported along D4:
        \begin{align*}
          &\forall s.\;\forall p\in P(s).\;\forall q\in P(s\sub p).\;
          \forall q'\in P\bigl((s\sub p)\sub q\bigr).\\
          &\qquad
          (p\shift q)\shift\bigl((\mathrm{D4}_{s,p,q})^{-1}\cdot q'\bigr)
          \;=\;
          p\shift\bigl(q\shift q'\bigr).
        \end{align*}
\end{enumerate}
\end{definition}

\begin{remark}
Laws D1 and D4 are equalities of \emph{shapes}. Laws D2 and D5 equate
\emph{positions} living in fibres that are identified only up to those shape
equalities, so the statements necessarily involve transport. Law D3 needs no
transport: both sides inhabit $P(s)$ directly. The numbering D1--D5 follows
the literature; the logical dependencies are that D2 refers to D1 and D5
refers to D4.
\end{remark}

\begin{example}
\begin{enumerate}
  \item \textbf{Non-empty lists / the store comonad on finite positions.}
        Shapes are non-zero natural numbers; positions of $n$ are
        $\{0,\ldots,n-1\}$; the root is $0$; the sub-shape under $p$ is
        $n-p$; the shift is addition (truncated to the remaining length).
  \item \textbf{Streams.} A single shape, positions $\mathbb{N}$, root $0$,
        sub-shape always the unique shape, shift ordinary addition.
  \item \textbf{Tree zippers.} Shapes are trees; positions are paths from
        the root; the root position is the empty path; sub-shape is the
        subtree at a path; shift is path concatenation.
\end{enumerate}
\end{example}

%======================================================================
\section{The induced comonad}
%======================================================================

Fix a directed container $D$ with underlying container $C$, and write
$\Ext\,C$ for its extension.

\begin{definition}[Counit]\label{def:counit}
The \textbf{counit} of the comonad is the family of maps
\begin{equation}
  \varepsilon_X \;:\; \Ext\,C\,X \;\longrightarrow\; X,
  \qquad
  \varepsilon_X\langle s,\,v\rangle \;=\; v\bigl(\rootpos(s)\bigr).
\end{equation}
It reads the payload at the root position of the shape.
\end{definition}

\begin{definition}[Comultiplication]\label{def:comult}
The \textbf{comultiplication} is the family of maps
\begin{equation}
  \delta_X \;:\;
  \Ext\,C\,X \;\longrightarrow\; \Ext\,C\bigl(\Ext\,C\,X\bigr)
\end{equation}
defined by
\begin{equation}
  \delta_X\langle s,\,v\rangle
  \;=\;
  \Bigl\langle
    s,\;
    p \;\mapsto\;
    \bigl\langle
      s\sub p,\;
      q \;\mapsto\; v(p\shift q)
    \bigr\rangle
  \Bigr\rangle.
\end{equation}
The outer shape is unchanged. At each position $p$ of $s$ one places the
element of $\Ext\,C\,X$ whose shape is the sub-shape under $p$ and whose
labelling is $v$ reindexed by shifting from $p$.
\end{definition}

%======================================================================
\section{Comonad laws}
%======================================================================

\begin{theorem}[Left counit]\label{thm:left-counit}
For every type $X$,
\begin{equation}
  \varepsilon_{\Ext\,C\,X}\;\circ\;\delta_X
  \;=\;
  \id_{\Ext\,C\,X}.
\end{equation}
\end{theorem}

\begin{proof}
By lean.
\end{proof}

\begin{remark}[Idea of the argument]
Evaluating the composite on $\langle s,v\rangle$ yields
\[
  \bigl\langle
    s\sub\rootpos(s),\;
    r\mapsto v\bigl(\rootpos(s)\shift r\bigr)
  \bigr\rangle.
\]
Law D1 equates the shape $s\sub\rootpos(s)$ with $s$. Law D2 equates
$\rootpos(s)\shift r$ with the transport of $r$ along D1, so the labelling
agrees with $v$ after transport. Lemma~\ref{lem:ext-eq} concludes.
\end{remark}

\begin{theorem}[Right counit]\label{thm:right-counit}
For every type $X$,
\begin{equation}
  \Ext\,C\,\varepsilon_X\;\circ\;\delta_X
  \;=\;
  \id_{\Ext\,C\,X}.
\end{equation}
\end{theorem}

\begin{proof}
By lean.
\end{proof}

\begin{remark}[Idea of the argument]
The composite sends $\langle s,v\rangle$ to
\[
  \bigl\langle s,\;
    q\mapsto v\bigl(q\shift\rootpos(s\sub q)\bigr)
  \bigr\rangle.
\]
Shapes are unchanged. Law D3 says $q\shift\rootpos(s\sub q)=q$, so the
labelling is $v$.
\end{remark}

\begin{theorem}[Coassociativity]\label{thm:coassoc}
For every type $X$,
\begin{equation}
  \Ext\,C\,\delta_X\;\circ\;\delta_X
  \;=\;
  \delta_{\Ext\,C\,X}\;\circ\;\delta_X.
\end{equation}
\end{theorem}

\begin{proof}
By lean.
\end{proof}

\begin{remark}[Idea of the argument]
Both sides produce, from $\langle s,v\rangle$, a thrice-nested element of
the extension. The outermost shapes agree. At an outer position $p$ and an
inner position $q$, the two sides place shapes
$(s\sub p)\sub q$ and $s\sub(p\shift q)$ respectively; these are identified
by D4 (up to symmetry). At the deepest level the labellings are
$v\bigl(p\shift(q\shift r)\bigr)$ and
$v\bigl((p\shift q)\shift r\bigr)$ (with $r$ transported along D4 on one
side); these agree by D5. Applying Lemma~\ref{lem:ext-eq} thrice yields the
result.
\end{remark}

\begin{corollary}[Comonad structure]\label{cor:comonad}
If $D$ is a directed container with underlying container $C$, then
$(\Ext\,C,\,\varepsilon,\,\delta)$ is a comonad on~$\Type$.
\end{corollary}

\begin{proof}
By lean. (The three laws are Theorems~\ref{thm:left-counit},
\ref{thm:right-counit}, and~\ref{thm:coassoc}.)
\end{proof}

%======================================================================
\section{Discussion}
%======================================================================

\paragraph{What has been established.}
Directed structure on a container is precisely the data needed to make the
extension a comonad. The correspondence is tight: each comonad law is the
direct image of one or two directed laws, and the only non-trivial
bookkeeping is transport of positions along shape equalities (isolated in
Lemma~\ref{lem:ext-eq}).

\paragraph{Dependence and transport.}
The laws D2 and D5 are the places where dependent type theory is essential:
positions of a sub-shape of a sub-shape are not definitionally positions of a
shifted sub-shape, so associativity of shift must be stated with an explicit
transport along D4. In set-theoretic presentations this transport is
invisible; in a formal development it is the content of the calculation.

\paragraph{Scope.}
This note treats a single directed container and the comonad on its
extension. Morphisms of directed containers, their identification with
cofunctors of the presented categories, and the monoidal structure on the
category of containers of which directed containers are the comonoids, belong
to companion developments.

\paragraph{Machine checking.}
Every definition and theorem statement above corresponds to a declaration in a
machine-checked formalisation. The phrase ``by lean'' indicates that the proof
has been verified by a proof assistant; we have not reproduced the formal
proof scripts here.

%======================================================================
\section*{Acknowledgements}
%======================================================================

The mathematical content follows the theory of directed containers of Ahman,
Chapman, and Uustalu. The presentation here is a prose rendering of a formal
development intended for human review of the theorem statements.

\begin{thebibliography}{9}

\bibitem{ahman2012}
D.~Ahman, J.~Chapman, and T.~Uustalu.
\newblock When is a container a comonad?
\newblock In \emph{FoSSaCS}, volume~7213 of \emph{LNCS}, pages 74--88.
\newblock Springer, 2012.

\bibitem{ahman2014}
D.~Ahman, J.~Chapman, and T.~Uustalu.
\newblock When is a container a comonad?
\newblock \emph{Logical Methods in Computer Science}, 10(3), 2014.

\bibitem{abbott2005}
M.~Abbott, T.~Altenkirch, and N.~Ghani.
\newblock Containers: constructing strictly positive types.
\newblock \emph{Theoretical Computer Science}, 342(1):3--27, 2005.

\end{thebibliography}

\end{document}
